%% APPENDIX B -- the spatial screen: calibration, validation, and the one
%% worked example of the vetting chain.  Owner: appendix-triage.
%% Carries \label{app:radius} (cited by 04_method.tex) and
%% \label{sec:heldout} (cited from the other appendices).
\section{The spatial screen: calibration and validation}
\label{app:radius}

The spatial screen asks whether an excess sits at the star or anywhere in the
field, and it removes most of the unattributed crossings. Everything the
dispositions rest on therefore depends on how well that screen is
calibrated, which is what this appendix measures. The answer is adverse to
the screen and is stated as such: \rankstatus{}

The star lies at the phase centre, inside the control annulus, and the
asymmetry is real. \emph{It is not a primary-beam effect.} Spectra are
extracted from phase-rotated visibilities in apparent flux and are never
corrected for the primary-beam response
(Appendix~\ref{app:conventions}), so the familiar rise of image noise with
distance from the pointing centre cannot operate here; and were it
operating, it would make the control probes noisier than the star and so
place the stellar rank \emph{above} one half, which is the opposite of what
is measured. \textbf{No mechanism for the displacement is established, and
we prefer to say so.} What the control ensemble does show is a weak rise of
its own level with radius, and the obvious repair is to remove it: each
control probe is standardised within its own radius bin, the probe radii
being known from the fixed draw, with no free parameter.

\emph{Applied out of sample, the correction does not repair the statistic.}
It was frozen --- probe radii, bin edges and standardisation alike --- and
applied unmodified to \RadValStarN{} windows that played no part in its
construction. The uncorrected stellar rank has median \RadValStarRawMed{}
against the exchangeable value of 0.5; after correction the median is
\RadValStarCorMed, moving \RadValStarToward{} it, and the result holds
separately in both out-of-sample partitions. Rescoring the whole survey with
the gradient removed moves the rank-passing list from \LocStageGlobal{}
windows to \LocStageLocal{} --- losing \LocStageLost{} and gaining
\LocStageGained, all marginal and all dispositioned on other grounds --- and
leaves the survey median rank slightly further from one half than before.

The interpretation is narrow and, we think, unavoidable. The radial gradient
is real, but removing it leaves the median no closer to one half, so the
residual is not radial, and neither normalisation yields a calibrated tail
probability. \textbf{We therefore retain the frozen screen as primary --- it
has at least the virtue of having been fixed in advance --- use it only to
decide which windows are examined, and rest every disposition on physical
evidence or on recurrence rather than on either rank.} The radial
correction, the region-maximum audit and the pseudo-star tests each leave
every disposition unchanged.

\emph{What the screen costs is not uniform, and almost all of it is paid in
a handful of fields.} Requiring the star to outrank its ring moves the
Class~A ninety-per-cent recovery point, taken as a single blanket factor,
from $\ScrTrigOnlyA\,P_{\rm trig}$ to $\EirpNinetyMultA\,P_{\rm trig}$, a
factor $\times\ScrCostA$. Where the ring
is noise the cost is only $\times\ScrCostNoise$, screening against a ring
peaking at $C$ times the noise being equivalent to raising the trigger to
$C$; almost the whole remaining effect comes from \ScrNBright{} windows of
\ScrNWinA{} whose rings carry resolved disc emission at
\ScrBrightLo--\ScrBrightHi$\sigma$ (\ScrBrightStars), where no carrier near
the trigger can pass at all.

\emph{One mechanism can be excluded by measurement, and it is the one that
would matter most.} Equation~\ref{eq:tstar} divides every position by a
single noise scale, the median absolute deviation across the control probes,
shared by the star and all of them. A star bright enough to carry
multiplicative calibration residuals of its own would then have a noise that
the shared scale does not describe, and its statistic would be inflated with
nothing in the search to reveal it. The ratio of the two noises is
measurable from the released control vectors alone: because the star and
each probe are divided by the same scale and maximised over the same
channel$\times$drift grid, the stellar statistic divided by the median of
the \NCtrl{} control statistics \emph{is} that ratio, and it is one when the
two noises agree. Over the \NsrNWin{} windows carrying no threshold crossing
its median is \NsrMed, with a star-clustered 95 per cent interval of
\NsrLo--\NsrHi. It does not climb with the brightness of the emission in the
field, which is what a multiplicative error would scale with: \NsrFifthLo{}
in the fifth of windows with the faintest control rings against
\NsrFifthHi{} in the brightest fifth, and \NsrBrightMed{} over the
\NsrBrightN{} windows whose ring exceeds \NsrBrightSig$\sigma$ and is
therefore set by resolved emission rather than by noise. That the ratio can
rise at all is shown by the \NsrCrossN{} windows that do carry a crossing,
where it is \NsrCrossMed. Part of the reason it does not rise otherwise is
structural: a gain error multiplying a continuum raises every channel of an
integration by the same factor, and the per-integration median along
frequency removes exactly that. \textbf{A stellar-position noise denominator
would therefore move no statistic in this survey, and the search has not
been re-run with one.}

\subsection{The reserved hold-out, and what it calibrates}
\label{sec:heldout}

The statistic and the mask were settled on the searched sample and cannot
also calibrate themselves, so a hold-out was reserved by a rule fixed and
recorded before the archival search finished and before any reserved block
was searched: a block is held out if and only if
$\mathrm{sha256}(\textsc{uid})_{[:8]} \bmod 5 = 0$, subject to two guards
that use identifiers alone --- a block on the pre-registered block list
deposited with this paper stays in the survey, and no system loses every
block --- so that the assignment can be recomputed from the UIDs and that
list alone. The three choices that can manufacture a candidate were fixed
before the rule itself: the statistic of Eq.~\ref{eq:tstar}, the promote and
retire criteria, and then the reservation, with the first reserved block
searched \HoLeadBlockDays{} days later. Every diagnostic built afterwards
--- the radial normalisation, the visibility-domain test, the drift-resolved
injection analysis, the $P_{90}$ budget --- is incapable of promoting a
window.

Of \NEbProcessedAll{} blocks processed, \NEbHeldOut{} are reserved
(\HoPctHeld{} per cent), giving \HoWin{} windows toward \HoStars{} stars with
no star and no system lost from the headline sample. Three quantities are
calibrated on those blocks and on nothing else.

First, the rank displacement. Out of sample the stellar rank has median
\HoRankMed{} where one half is expected, with a block-clustered 95 per cent
interval of \HoRankMedLo--\HoRankMedHi{} over \HoWin{} windows, and
\HOFracLoPct{} per cent of ranks below 0.1 against \HOFracHiPct{} per cent
above 0.9, so the distribution is displaced rather than growing a tail. The
displacement runs the way that matters: a rank below one half places the star
high, and the star in fact outranks all \NCtrl{} probes in \NsrFirstObs{} of
the survey's windows against \NsrFirstExp{} if the positions were
interchangeable. The direction therefore inflates star-first outcomes rather
than suppressing them, which is why the chance expectations of
Appendix~\ref{app:falsealarm} are calibrated on the measured rate.
Clustering removes neither the shift nor its size: \HOBlockBelow{} of
\HOBlockN{} blocks have a mean rank below 0.5, and a star-clustered bootstrap
places the median at \HOBootMed{} (95 per cent \HOBootLo--\HOBootHi). The
searched sample leans the same way without resolving it (median \SurvMedP,
star-clustered 95 per cent \SurvBootLo--\SurvBootHi), and the two
distributions are indistinguishable, so the in-sample checks were
underpowered rather than contradicted. A displacement that survives a
hold-out chosen before the data were seen is a property of the statistic and
not of this sample, which is the quantitative reason the rank is a screen and
never a significance.

Second, the false-alarm tail factor --- the measured rate at which a
no-signal position ranks first, divided by the rate exchangeability would
give. It is \HoTailFactor{} on \HoPseudoTrials{} trials, with a bootstrap
clustered over the \HoTailBlocks{} reserved blocks giving
\HoTailLo--\HoTailHi. That factor and its width set the accuracy of every
chance expectation in this paper. The reserved blocks contain \HoStageOne{}
rank-passing crossings against \HoStageOneExp{} expected.

Third, the radial behaviour of the control ensemble. The standardised
control level runs \HoRadBins{} from the inner to the outer edge of the
annulus, a mean per-window slope of $\HoRadSlope\pm\HoRadSlopeSe$: the
displacement reproduces out of sample, but the steep inner-edge excess seen
in sample does not, so the radial term is small, of the wrong sign to account
for the displacement, and removing it leaves the displacement behind.

A second out-of-sample set is used here and must not be confused with the
hold-out. The \emph{external calibration sample} is \CampBlocks{} execution
blocks toward \CampStars{} stars, reserved by a rule fixed before the
archival search finished; it reaches stars outside the 40\,pc work list, so
it enters no count of the science sample, and it exists only to calibrate the
rank statistic and to supply repeat observations the hold-out does not have.
Over its \CampWindows{} windows the frozen pipeline returns \CampStageOne{}
rank-passing crossings --- \CampStageOneCO{} of them $\beta$~Pictoris CO,
the positive control recovering on data it had never seen --- and
\CampUnattrib{} unattributed against \CampExpUn{} expected at the measured
tail rate.

\subsection{One crossing through the whole chain}
%% (subsection label retired: nothing cites it)

CP$-$72~2713 is the clearest worked example of the chain the survey
contains, because every stage returns a verdict rather than a shrug. In
Band~7 the crossing reaches $T_\star=\VnCpT$ in channel \VnCpChan{} of its
window and outranks every one of the \NCtrl{} control positions, so it is a
candidate; no transition of the frozen mask lies within
$\pm\MaskHalfKms$\,km\,s$^{-1}$ of it, so it is unattributed; and the
drift-following visibility fit puts the emission at the stellar position,
${\rm Re}/\sigma=\VnCpReAdopted$ against a control-position maximum of
\VnCpCtrlAdopted, which by the argument of \S\ref{sec:vistest} is a
consistency check and not a promotion.

\emph{What disposes of it is the repeat block}, and the number quoted for
that needs its method attached. The second block matches the first in tuning,
channel width, channel count, integration count and on-source time,
registers to \CpRecOffChan{} of a channel, and is the deeper of the two, so a
persistent emitter at the first block's flux would have returned
$T_\star=\CpRecExpT$ where it returns \CpRecT. Read as two measurements of
one constant amplitude --- the conservative reading, and the one this paper
quotes --- the pair is inconsistent at $\CpRecExclPair\sigma$; taking the
discovery flux as exact instead would give $\CpRecExclCond\sigma$. The blocks
are \CpSepH\,h apart, so this is a within-night test: it excludes a
persistent emitter of that strength on that baseline and nothing about an
intermittent one. \textbf{The disposition rests on the recurrence test and
not on the rank}, which is the chain applied to every crossing this survey
cannot attribute.

%% ------------------------------------------------------------------
%% REQUEST (appendix-triage -> results, 05_results.tex):
%%   CP-72 2713 appears nowhere in the body.  It is one of the two
%%   unattributed crossings that PASS the spatial screen, and the ledger's
%%   own chain string for it is "unattributed; rank-flagged; tested and does
%%   not recur".  It is quoted here as a worked example only.  If the results
%%   section means to name the two rank-passing crossings -- and it should --
%%   this is one of them, and the recurrence exclusion at \CpRecExclPair
%%   sigma is the result, not the rank.
%% NOTE (pre-existing contradiction, now fixed here):
%%   this paragraph previously gave T_star = 5.81 against "a ring maximum of
%%   5.68" and then concluded the crossing "still fails the rank screen".
%%   5.81 < 5.68 is false and so was the conclusion: ledger_v403.csv has
%%   rank_screen = True for this row.  The ring maximum had no macro, so the
%%   statement is now made in the form the screen actually tests -- the star
%%   outranks all \NCtrl controls -- and the six typed literals in the
%%   paragraph (5.81, 5.68, the SO offset, the transitions per GHz, the
%%   91 per cent coincidence rate and the 3-against-2 channel width) are
%%   gone: \VnCpT replaces the first, and the sentences carrying the others
%%   are deleted because no macro exists for them.
%% REQUEST (appendix-triage -> numbers):
%%   If the SO line-density argument is wanted back, it needs three macros:
%%   transitions per GHz over this window, the coincidence probability at
%%   +/- \MaskHalfKms km/s, and the feature's recovered width in channels
%%   against the instrumental response.  Without them it cannot be typeset.
%% ------------------------------------------------------------------
